\chapter{Diskusija}
U prethodnim poglavljima predstavljen je koncept dizajna proizvoda ili usluga sa jedinstvenim identifikatorom u okviru sistema elektronske trgovine, zasnovan na primeni NFT standarda na Eterijum blokčejnu kao najzastupljenijoj platformi za razvoj decentralizovanih aplikacija. Poseban akcenat stavljen je na optimizaciju operacija nad skladištem podataka, pri čemu se model mape vlasništva NFT standarda proširuje enkodiranjem dodatnih informacija o proizvodima ili uslugama u preostali prostor skladišne jedinice, sa ciljem minimizacije računarskog napora potrebnog za izvršavanje korisničkih transakcija.

S obzirom na to da Eterijum predstavlja javnu blokčejn mrežu sa visokim nivoom decentralizacije, koja obezbeđuje povećanu otpornost sistema i distribuiranu kontrolu, neminovno se javlja problem transakcionih troškova za krajnje korisnike. Ovaj problem je u radu detaljno analiziran, pri čemu su predložene odgovarajuće metode optimizacije sa ciljem smanjenja troškova transakcionih naknada, čime se direktno unapređuje korisničko iskustvo.

U ovom poglavlju razmatraju se mogućnosti za dalje unapređenje i proširenje funkcionalnosti pametnih ugovora u okviru NFT-zasnovane platforme za elektronsku trgovinu. Diskusija obuhvata dodatne funkcionalnosti NFT standarda koje nisu od primarnog značaja, kao i primenu postojećih, testiranih i optimizovanih modula. Takođe se analiziraju obrasci modularnog dizajna pametnih ugovora u kontekstu razvoja kompleksnih decentralizovanih aplikacija koje treba da budu dugoročno održive, kao i integracija sa spoljnim izvorima podataka, imajući u vidu da blokčejn sistemi inherentno funkcionišu kao zatvoreni sistemi.

Na kraju se analiziraju ograničenja predložene metodologije obrađene u prethodnom poglavlju, kao i dodatna razmatranja proistekla iz ove diskusije.

\section{Proširene funkcionalnosti NFT standarda u sistemima e-trgovine}
NFT standardi takođe predviđaju mehanizam tzv. zakačaljki \eng{hooks}, odnosno funkcija koje se pozivaju pre i nakon svake osnovne operacije, kroz interne funkcije \textit{\_beforeTokenTransfer} i \textit{\_afterTokenTransfer}. Ove funkcije su podrazumevano prazne, ali omogućavaju proširenje osnovnog ponašanja bez izmene standardne logike. Na primer, pre izvršenja prenosa može se implementirati dodatna provera autorizacije, kojom se potvrđuje da pošiljalac ili primalac pripada unapred definisanoj grupi korisnika.

Pored osnovnih operacija kreiranja, prenosa i uništavanja, NFT standardi predviđaju i funkciju odobravanja \eng{approve}, koja omogućava vlasniku NFT-a da ovlasti drugu adresu da prenese konkretan token u njegovo ime. Time se uvodi mehanizam imenovanja prava raspolaganja, bez prenosa vlasništva. Zbog prilagodljive prirode pametnih ugovora, ova funkcionalnost može se proširiti ili modifikovati, na primer tako da ovlašćena adresa može izvršiti i operaciju uništavanja tokena u ime vlasnika, u skladu sa poslovnom logikom platforme.

NFT aplikacije se oslanjaju na mehanizam metapodataka, pri čemu se dodatne informacije o tokenu čuvaju na eksternim servisima u JSON formatu, pri čemu ti podaci nisu direktno dostupni unutar same blokčejn mreže. Pristup tim podacima za određeni identifikator ostvaruje se putem funkcije \textit{tokenURI(uint256 tokenId)}, koja vraća referencu (najčešće URL) ka resursu dostupnom van mreže \cite{salem2025hidden}. U zavisnosti od arhitekture sistema, servisi za skladištenje mogu biti centralizovani (npr. baze podataka) ili decentralizovani (npr. IPFS).

Budući da se transakcione naknade plaćaju u primarnoj kriptovaluti blokčejn mreže, korisnik se može naći u situaciji da poseduje NFT, ali nema dovoljno sredstava za pokrivanje troškova izvršavanja transakcije. Ovaj problem moguće je ublažiti fleksibilnijim modelom interakcije primenom ekstenzije za \textit{dozvolu} \cite{NFTPERMITeip4494} \eng{permit}, koja je još u fazi standardizacije. Navedeni mehanizam omogućava vlasniku tokena da van mreže potpiše kriptografsku poruku kojom odobrava određenu operaciju, prosledi je drugoj strani, koja na osnovu tog potpisa može izvršiti odobrenu operaciju u njegovo ime i snositi troškove transakcione naknade.

Dodatno, ukoliko je NFT tržišno prenosiv, moguće je u sam pametni ugovor ugraditi mehanizam kojim se pri svakoj sekundarnoj prodaji ili pretprodaji automatski alocira određeni procenat sredstava ka unapred definisanoj adresi \cite{MARKETeip2981}. Ovakav pristup omogućava implementaciju modela trajnih autorskih naknada \eng{royalties}, čime se obezbeđuje kontinuirano finansiranje izdavaoca proizvoda. Time se finansijska logika distribucije prihoda ugrađuje direktno u izvršni kod pametnog ugovora, čime se eliminiše potreba za posrednicima i povećava transparentnost obračuna.

Za razliku od prethodno opisanih proširenja koja su specifično vezana za NFT standard, u praksi se široko primenjuju i generičke ekstenzije pametnih ugovora koje unapređuju bezbednost, kontrolu pristupa i upravljanje životnim ciklusom sistema. Ove komponente nisu ograničene na NFT logiku, već predstavljaju opšte obrasce modularnog dizajna primenljive na različite tipove pametnih ugovora. Biblioteka \textit{OpenZeppelin} obezbeđuje skup standardizovanih, bezbednosno proverenih i optimizovanih implementacija takvih ekstenzija, koje su postale industrijski standard.

Najčešće korišćene ekstenzije ove vrste uključuju:
\begin{enumerate}
\item \textit{Pausable} - mehanizam koji omogućava privremenu obustavu izvršavanja odabranih funkcija ugovora u vanrednim okolnostima, bez trajne izmene stanja sistema.
\item \textit{AccessControl} - model kontrole pristupa zasnovan na ulogama, kojim se precizno definišu privilegije i nadležnosti različitih uloga unutar sistema.
\item \textit{ReentrancyGuard} - zaštitni mehanizam koji sprečava ponovni ulazak u funkciju pre završetka njenog prethodnog izvršavanja, čime se sprečavaju napadi zasnovani na rekurzivnim pozivima \cite{samreen2020reentrancy}.
\item \textit{MultiSigner} - obrazac kolektivnog odlučivanja u kojem je za izvršavanje osetljivih operacija potrebno više nezavisnih ovlašćenih potpisa tzv. višepotpisnih novčanika, čime se smanjuje rizik centralizovanog upravljanja.
\end{enumerate}

Pored toga, moguće je delimično ili u određenim slučajevima u potpunosti ograničiti transparentnost javnih blokčejn mreža primenom potpuno homomorfnog šifrovanja \eng{fully homomorphic encryption}, koje omogućava izvršavanje operacija direktno nad šifrovanim podacima bez potrebe za njihovim dešifrovanjem \cite{fhevm}. U ovom modelu pametni ugovori primaju šifrovane ulaze, izvršavaju operacije nad njima i održavaju šifrovano stanje, bez pristupa originalnim vrednostima, koje mogu biti selektivno dekriptovane od strane ovlašćenih aktera. Tokom obrade transakcija u okviru konsenzusnog mehanizma operacije se izvršavaju nad šifrovanim vrednostima, pri čemu svi čvorovi deterministički dolaze do istog rezultata. Ovakav pristup predstavlja značajan korak ka širem usvajanju blokčejn tehnologije, iako se i dalje nalazi u fazi aktivnog razvoja.

Korišćenjem ovih mehanizama pametnih ugovora, distribuirana platforma za e-trgovinu može obezbediti širok spektar funkcionalnosti, uz zadržavanje formalne konzistentnosti i sigurnosnih garancija javne blokčejn mreže gde se poslovna pravila mogu implementirati direktno na nivou ugovora.

\section{Obrasci modularnog dizajna pametnih ugovora}
U arhitekturi platforme moguće je predvideti postojanje više međusobno povezanih pametnih ugovora, pri čemu pojedini ugovori preuzimaju specifične sistemske uloge. Ugovor o registru \cite{xu2018pattern} \eng{registry contract} predstavlja pametni ugovor koji funkcioniše kao decentralizovana baza podataka, čuvajući mape između imena, identifikatora ili drugih podataka i odgovarajućih adresa pametnih ugovora. Na taj način omogućava centralizovano usmeravanje poziva ka postojećim NFT kolekcijama.

Ukoliko je potrebno masovno kreiranje novih instanci kojima upravlja ugovor o registru, mogu se primeniti sledeći obrasci:
\begin{enumerate}
\item \textbf{Fabrički ugovor} \eng{factory contract} je pametni ugovor čija je osnovna funkcija postavljanje novih ugovora na mrežu. Njegova prednost ogleda se u mogućnosti dinamičkog kreiranja proizvoljnih instanci po potrebi, pri čemu je za postavljanje dovoljan bajtkod ugovora koji se implementira.

\item \textbf{Minimalni zastupnički ugovor} \cite{CLONEeip1167} \eng{minimal proxy contract} predstavlja specifičan slučaj kada se kreiraju identične instance ugovora, poznate kao klonovi \eng{clones}. U ovom obrascu izvršna logika nalazi se u jedinstvenom implementacionom ugovoru \eng{implementation contract}, dok klonovi poseduju sopstveno skladište podataka. Izvršni kod se delegira ka implementacionom ugovoru, čija se adresa nalazi u bajtkodu klona. Budući da je izvršni deo bajtkoda minimalnog zastupničkog ugovora dugačak svega 45 bajtova, troškovi postavljanja značajno su niži u odnosu na potpuno postavljanje ugovora sa identičnom logikom.
\end{enumerate}

Iako je kod pametnih ugovora po svojoj prirodi neizmenljiv nakon njihovog postavljanja na mrežu, ovo ograničenje može se prevazići primenom nadogradivih arhitektura \cite{meisami2023comprehensive}. Uobičajeni pristup podrazumeva razdvajanje skladišne i izvršne logike kroz dva međusobno povezana ugovora:

\begin{enumerate}
\item \textbf{Zastupnički ugovor} \cite{PROXYeip1967}, koji čuva skladišne promenljive i predstavlja stabilnu tačku interakcije sa korisnicima.
\item \textbf{Implementacioni ugovor}, koji sadrži izvršni kod i može se zameniti u procesu nadogradnje bez promene adrese zastupničkog ugovora.
\end{enumerate}
U zavisnosti od načina implementacije nadogradnje razlikuju se transparentni zastupnički ugovor \eng{Transparent Proxy} i univerzalni nadogradivi zastupnički ugovor \eng{Universal Upgradeable Proxy Standard, UUPS}.

Kod kompleksnijih sistema mogu se primeniti dodatni obrasci:
\begin{enumerate}
\item U slučaju postojanja većeg broja zastupničkih ugovora, radi pojednostavljenja nadogradnje može se primeniti obrazac svetionika \eng{beacon}. U tom slučaju zastupnički ugovori u svom bajtkodu čuvaju adresu ugovora svetionika, koji centralno upravlja adresom aktuelne implementacije. Izmena implementacije na nivou ugovora svetionika automatski utiče na sve povezane zastupničke ugovore.

\item Kada je potrebno modularizovati implementacioni ugovor i prevazići ograničenje maksimalne veličine bajtkoda, može se primeniti dijamantski obrazac \cite{DIAMONDeip2535} \eng{diamond pattern}. Ovaj pristup omogućava razdvajanje funkcionalnosti na više međusobno povezanih komponenti \eng{facets}, pri čemu centralni ugovor \eng{diamond} delegira pozive odgovarajućim modulima. Time se omogućava nadogradnja pojedinačnih delova sistema i fleksibilno proširenje funkcionalnosti.
\end{enumerate}

U praksi, proces nadogradnje i upravljanja implementacijama često se organizuje kroz višepotpisnih novčanika ili kroz DAO mehanizme, kako bi se smanjila centralizovana kontrola i povećala decentralizacija odluka o promenama koda, posebno u slučajevima kada promena utiče na velike korisničke baze.

Ovakav arhitektonski pristup omogućava da pametni ugovori preuzmu ulogu aplikativnog sloja platforme, pri čemu korisnici, nakon autentikacije putem kriptografskih ključeva, neposredno komuniciraju sa blokčejn infrastrukturom, dok se logika sistema izvršava u determinističkom okruženju distribuirane mreže.

\section{Integracija blokčejn sistema sa spoljašnjim izvorima podataka}
Blokčejn sistem predstavlja zatvoreno determinističko okruženje u kome svaki čvor, nakon prijema novog bloka, mora izračunati identično globalno stanje sistema. Ovakav model obezbeđuje visok stepen integriteta i neizmenljivosti podataka, jer naknadna manipulacija već potvrđenim zapisima praktično nije moguća. Međutim, posledica takve arhitekture jeste \textit{odsustvo ugrađenih mehanizama za pristup informacijama van same mreže}, što značajno ograničava funkcionalni domet pametnih ugovora.

Da bi se podaci javno dostupni na klasičnom internetu uveli u blokčejn sistem i učinili dostupnim pametnim ugovorima, koristi se mehanizam tzv. proročišta \cite{beniiche2020study} \eng{oracles}, koji predstavlja posrednički sloj između blokčejn sistema i spoljašnjih izvora podataka. S obzirom na to da blokčejn ne može direktno pozivati eksterne API servise, proročište tu operaciju izvršava van mreže, preuzima relevantne podatke iz eksternih izvora, obrađuje ih i potom putem transakcije upisuje rezultat u sopstveni pametni ugovor na blokčejn mreži. Ostali pametni ugovori zatim pristupaju tim podacima pozivanjem funkcija ugovora proročišta.

Osnovni izazov ovakvog modela jeste obezbeđivanje poverenja i očuvanje decentralizacije. Ukoliko decentralizovani blokčejn sistem zavisi od centralizovanog proročišta, uvodi se jedinstvena tačka otkaza, što je u literaturi poznato kao \textit{problem proročišta} \eng{The Oracle Problem}. Kompromitovano proročište ili nepouzdan izvor podataka može uticati na izvršavanje svih pametnih ugovora koji se na njega oslanjaju, što potencijalno dovodi do bezbednosnih incidenata i finansijskih gubitaka \cite{arora2024secplf}.

Radi prevazilaženja ovog problema razvijene su decentralizovane mreže proročišta \cite{breidenbach2021chainlink} \eng{Decentralized Oracle Networks, DON}, koje agregiraju podatke iz više nezavisnih čvorova i različitih izvora. Agregacijom i mehanizmima konsenzusa umanjuje se rizik manipulacije pojedinačnog entiteta, čime se povećavaju pouzdanost i bezbednost sistema.

Pored dostavljanja eksternih podataka, decentralizovane mreže proročišta omogućavaju i druge funkcionalnosti, uključujući generisanje verifikabilne slučajnosti, automatizaciju izvršavanja pametnih ugovora, međulančanu komunikaciju, izvršavanje određenih proračuna van lanca, kao i mehanizme dokaza o rezervama. Time se značajno proširuje primenljivost blokčejn sistema u realnim poslovnim okruženjima.

Mehanizmi proročišta su danas visoko razvijeni kada je reč o javno dostupnim podacima. Međutim, u savremenim primenama javlja se potreba za pouzdanim uvozom privatnih korisničkih podataka iz postojećih informacionih sistema, kao što su bankarski servisi, sistemi elektronskog identiteta, društvene mreže, servisi za elektronsku trgovinu i drugi centralizovani internet servisi.

U okviru standardnog TLS protokola, nakon uspostavljanja sigurne komunikacione sesije, podaci koje klijent preuzima od servera jesu kriptografski zaštićeni tokom prenosa, ali nisu kriptografski dokazivi trećoj strani. Drugim rečima, klijent ne poseduje dokaz koji bi nezavisno potvrdio da su konkretni podaci zaista potekli od određenog servera.

Nadogradnja TLS protokola, poznata kao zkTLS \cite{kireieva2024zktls}, uvodi mehanizam generisanja kriptografskog dokaza kojim se potvrđuje da su pribavljeni podaci zaista dostavljeni od strane servera koji poseduje važeći digitalni sertifikat. Ovako definisan pristup omogućava da se selektivno izdvojene informacije iz postojećih internet servisa, koji već koriste TLS infrastrukturu, pouzdano prenesu u blokčejn sistem zajedno sa kriptografskim dokazom o njihovom poreklu, bez potrebe za modifikacijom samih servisa. Na taj način ostvaruje se visok potencijal za globalno skaliranje blokčejn sistema.

Ipak, bezbednosni model ostaje uslovljen pretpostavkom ispravnosti i integriteta eksternih servisa koji generišu podatke, kao i poverenjem u autoritete za sertifikate u okviru TLS infrastrukture. Stoga se sigurnost sistema delimično oslanja na spoljne entitete, što predstavlja dodatnu tačku rizika u ukupnoj arhitekturi.

\section{Decentralizovano skladištenje podataka}
Za razliku od tradicionalnih web aplikacija, kod kojih se sadržaj preuzima sa unapred poznatog servera, IPFS koristi distribuiranu mrežu čvorova u kojoj se podaci pronalaze na osnovu jedinstvenog identifikatora sadržaja (CID). Takav pristup doprinosi otpornosti sistema i eliminiše zavisnost od centralnog servera, ali istovremeno uvodi određene izazove koji mogu uticati na korisničko iskustvo i praktičnu primenu.

Jedan od najvažnijih izazova odnosi se na vreme odaziva prilikom pristupa sadržaju. Za razliku od centralizovanih mreža za isporuku sadržaja (CDN), kod kojih se sadržaj najčešće isporučuje sa geografski najbližeg servera, IPFS prvo mora pronaći odgovarajuće čvorove, nakon čega započinje samo preuzimanje podataka. Vreme potrebno za preuzimanje datoteke zavisi od broja čvorova koji poseduju traženi sadržaj, njihove dostupnosti, kao i mrežne udaljenosti između korisnika i odgovarajućih čvorova. Ukoliko se određeni sadržaj nalazi na malom broju čvorova, vreme odziva može biti značajno veće.

Iako ovakvo ponašanje ne utiče na funkcionalnu ispravnost sistema, u elektronskoj trgovini čak i relativno mala povećanja vremena učitavanja mogu negativno uticati na korisničko iskustvo, imajući u vidu da korisnici očekuju brz odziv nezavisno od tehnologije na kojoj je aplikacija implementirana.

Jedan od pristupa za ublažavanje ovog problema jeste uvođenje keširanja na pristupnicima, čime se može značajno smanjiti vreme pristupa sadržaju. Merenja pokazuju da se na ovaj način približno 76\% zahteva može obraditi za manje od 250 ms \cite{trautwein2022design}.

Slično tome, određeno vreme je potrebno i da mreža utvrdi da traženi sadržaj ne postoji ni na jednom čvoru jer postoji određeno vreme potrebno za završetak procesa pretrage. Upit mora biti prosleđen između čvorova sve dok se ne pronađe odgovarajući sadržaj ili dok mreža ne zaključi da nijedan poznati čvor ne poseduje traženi CID, što predstavlja dodatni izvor vremena odaziva u odnosu na klasične centralizovane sisteme.

Dodatni izazov predstavlja trajnost podataka \eng{data persistence}, jer IPFS ne garantuje da će jednom postavljena datoteka biti trajno dostupna, već njena dostupnost zavisi od toga da li je neki čvor i dalje čuva u svom lokalnom skladištu. Oslanjanje na usluge kačenja uvodi određeni stepen centralizacije, budući da dostupnost podataka postaje zavisna od treće strane \cite{lamichhane2024verifiable}.

Zbog otvorene prirode IPFS mreže i distribuiranog rutiranja zahteva, zlonamerni učesnici mogu pokušati da prate koji sadržaj korisnici pretražuju ili preuzimaju, što utiče na privatnost korisnika tokom pristupa sadržaju. Savremeni pristupi unapređenju privatnosti zasnivaju se na prikrivanju identiteta korisnika i/ili sadržaja tokom procesa pronalaženja podataka, čime se otežava povezivanje zahteva sa konkretnim učesnicima mreže \cite{katsantas2024enhancing}, bez oslanjanja na pouzdane treće strane.

Sa aspekta privatnosti podataka, IPFS obezbeđuje integritet sadržaja putem kriptografskog heširanja, ali ne garantuje njegovu poverljivost. Svako ko poseduje odgovarajući CID može preuzeti sadržaj ukoliko je on dostupan u mreži. Iz tog razloga IPFS nije pogodan za skladištenje osetljivih podataka u izvornom obliku, već ih treba šifrovati pre postavljanja \cite{manakhari2024enhancing}, a ključ za dešifrovanje čuvati odvojeno. Na ovaj način se zadržavaju prednosti sadržajno adresiranog skladištenja, uz istovremeno očuvanje poverljivosti podataka.

U tradicionalnim sistemima datoteke se mogu organizovati prema nazivima, direktorijumima ili metapodacima, što omogućava jednostavnu pretragu sadržaja na osnovu sličnih naziva ili drugih atributa. Sadržajno adresirani sistemi ne obezbeđuju mehanizme za semantičku pretragu sadržaja, zbog čega decentralizovane aplikacije moraju koristiti dodatne baze metapodataka ili indekse koji mapiraju korisnički razumljive identifikatore na odgovarajuće CID vrednosti \cite{shevarev2025dynamic}.

Većina savremenih web pregledača ne podržava urođeni pristup IPFS sadržaju, zbog čega se u praksi često koriste pristupnici koji predstavljaju posrednika između korisnika i IPFS mreže. Slično tome, rezolucija ENS domena nije univerzalno podržana u svim pregledačima, pa je za pristup decentralizovanim aplikacijama često neophodno koristiti dodatne servise. Iako navedena ograničenja ne predstavljaju fundamentalne prepreke za funkcionisanje decentralizovanih aplikacija, ona pokazuju da potpuna decentralizacija prezentacionog sloja još uvek nije u potpunosti podržana od strane postojeće web infrastrukture.